\section*{Preface}

These notes cover the Analysis 1 course taught during Aug-Nov 2026 at CMI by Prof. Karandikar.
Throughout the course the emphasis will be on rigorous proofs of the fundamental theorems used in Calculus.
Construction of reals, and metric spaces will \emph{not} be covered.

\subsection{Timings}
$2:00-3:15$pm Tuesday, Thursday and $11:50-12:05$pm on Wednesday. Shared tutorials with Algebra every Monday and
Friday at $2$pm.

\subsection{Textbook}
\emph{Understanding Analysis} by Stephen Abbott is a good reference for this course. He recommends
\emph{Principles of Mathematical Analysis} by Walter Rudin, which I don't recommend as a first text.

\subsection{Syllabus}

\begin{enumerate}
\item Axioms of the real number system without construction, applications of
the least-upper-bound property, Archimedean principle, existence of $n$th roots of
positive real numbers, $a^x$ for $a > 0$ and $x > 0$, cardinality, countability of rational
numbers, uncountability of real numbers.
\item Convergence of sequences, uniqueness of limit, Sandwich lemma, exam-
ples, monotonic sequences, subsequences, Bolzano-Weierstrass theorem, lim sup and
lim inf, Cauchy sequences, completeness of $\mathbb{R}$.
\item Infinite series, convergence and divergence, absolute convergence, condi-
tional convergence, convergence/divergence of $\sum_{k=1}^{\infty} 1/k^p$, comparison test, root test,
ratio test, Leibniz test.
\item Complex numbers, power series, radius of convergence of power series.
\item Continuous functions on intervals of $\mathbb{R}$, intermediate value theorem, bound-
edness of continuous functions on closed and bounded intervals. Uniform continuity,
uniform convergence, uniform limit of continuous functions is continuous, counter
example for point wise convergence.
\item Differentiation, mean value theorem, Taylor's theorem, application of Taylor's theorem to maxima and minima,
L'Hôpital rule.
\item Construction of ez using power series, proof of the periodicity of sin and
cos.
\item Riemann Integration: Riemann integrals, Riemann integrability of continu-
ous functions, fundamental theorem of calculus. Improper integrals
\end{enumerate}

\subsection{Grading}

Homework will \emph{not} be graded. The grading will be distributed between quizzes, midsem and finals.

\subsubsection{Goal}

These notes are intended for an audience that is familiar with basic highschool mathematics,
in particular the notions of functions, basic set theory, etc. Beyond this, I try not
to assume any prior knowledge. The goal is to produce a more comprehensive, if slightly verbose,
introduction to analysis, that tries not to omit details that otherwise might be left out, and
tries to motivate the ideas presented through intuitive notions and visualizations.

For this reason, most exercises have a solution provided. I usually at least provide a sketch,
unless the exercise is very simple. I recommend that you try to solve the exercise yourself
for a while before looking at the solution. I have taken some liberty in rearranging the
questions from the course, so that some maybe a part of the main text itself, while some
are provided as exercises.

In particular, I hope this will serve as a reference for any of my classmates taking this course,
who have no prior background in analysis. If you find any errors, or have any
suggestions please mail me at \href{mailto:lunalgia@pm.me}{lunalgia@pm.me}.
